\section{The concrete encoding}
\label{sec:encode}

\texttt{encode.thy} instantiates the interface of \S\ref{sec:encoder-interface}
with a G\"odel numbering built from the Cantor pairing of \S\ref{sec:cpair},
supplies concrete recursive definitions for every function the checker is built
from, discharges all the assumptions, and instantiates the consistency theorem.

\subsection{The encoders}

A code is a pair whose first component is the constructor tag and whose second
is the payload.  Formulas use the pairing directly, and terms compose it with
an involution that exchanges $0$ and $1$.

\thy{enc-encoders}

\begin{remarkx}[Why \isa{swap01}]
The interface requires \isa{tag\_T 0 = T\_ZERO}, so the code $0$ must denote the
term $0$.  With \isa{T\_ZERO = 1} and \isa{cpx 0 = 0}, taking
\isa{tag\_T = cpx} would make the code $0$ a variable reference, since
\isa{T\_VAR = 0}, and the structural-decrease assumption
\isa{t ≠ 0 ⟹ load\_T t < t = 1} would then have no atomic base to stop at.
Composing with \isa{swap01} on both the packing and the unpacking side relabels
those two tags and leaves everything else alone, and
\isa{swap01 (swap01 t) = t} makes the composition cancel.  Formulas need no
adjustment, since \isa{F\_EQ = 0} already.
\end{remarkx}

\thy{enc-swap01}

\noindent
\isa{swap01\_N} and \isa{swap01\_inv} are proved by \isa{cases bool:} on the
guards.  The habeas quid premise \isa{t N} is what makes \isa{t = 0} grounded
and the case split available.

\subsection{Concrete definitions}

Every function the locales left abstract is introduced in a single
\isa{axiomatization} block whose axioms are their recursive \isa{:=} equations.
This is the definitional mechanism of \S\ref{sec:defmech} applied at scale, and
it is why a self-referential evaluator can be written down at all: no
termination obligation is incurred when the definition is made.

\thy{enc-signature}

\noindent
Thirty-three constants: the two evaluators, the three substitutions, the
freshness predicates, the assignment update, and the checker hierarchy down to
\isa{is\_valid\_proof}.  Their equations repeat the locale assumptions
verbatim, for instance the direct evaluator

\thy{enc-eval}

\noindent
and the assignment update

\thy{enc-asn-put}

\noindent
The background definition list stays opaque.  Only its termination is assumed:

\begin{center}
\isa{axiomatization dfns :: "dfn" where dfns\_N: "dfns N"}
\end{center}

\noindent
so the result holds for every terminating definition list, including lists
whose definitions diverge on all arguments.

\subsection{Discharging the assumptions}

What remains is the arithmetic that the locale structure confined to this one
file.

\paragraph{Habeas quid.}
Each of the six encoder operations preserves termination, by the corresponding
pairing lemma.

\thy{enc-hq}

\paragraph{Round-trips.}
Injectivity is the projection equations of \S\ref{sec:cpair} composed with
\isa{swap01\_inv}.  The retraction property is \isa{cpair\_reconstr}, that a
number is the pair of its projections.

\thy{enc-roundtrip}

\paragraph{Structural decrease.}
Both payload extractors are \isa{cpy}, so the assumption reduces to strict
decrease of the second projection on nonzero arguments, which is already
available in \texttt{GD.thy}.  The base case is
\isa{swap01 (cpx 0) = swap01 0 = 1}.

\thy{enc-decrease}

\paragraph{Monotonicity.}
This is the assumption that needs work.  The chain begins with
\isa{pair\_step}, that a single successor step in the second component does not
decrease the pair,
\begin{center}
\isa{pair\_step: ⟦a N; k N⟧ ⟹ ⟨a, k⟩ ≤ ⟨a, S k⟩ = 1},
\end{center}
which follows from \isa{cpair\_suc} and repeated use of
\isa{leq\_monotone\_add\_r}.  \isa{pair\_mono\_2} closes it under induction on
the second component, and the two interface obligations follow.

\thy{enc-mono}

\noindent
The induction in \isa{pair\_mono\_2} is a good miniature of the style of the
whole development.  The statement being inducted on is the object-level
implication \isa{(x ≤ y = 1) ⟶ (⟨a, x⟩ ≤ ⟨a, y⟩ = 1)}, every \isa{implI} step
first discharges a groundedness obligation through \isa{eqBool}, and the case
analysis in the step case goes through \isa{cases\_bool}.

\subsection{Instantiation and the theorem}

With the assumptions available the locale is interpreted.  The interpretation
supplies, in order, the six encoder operations, the three freshness predicates,
the definition list and its membership test, the step-indexed evaluator, the
three substitutions and the assignment update, and the checker hierarchy.

\thy{enc-interpretation}

\noindent
Every goal generated by \isa{unfold\_locales} is closed by \isa{fact}: the
structural obligations by the lemmas above, and the definitional obligations by
the \isa{:=} axioms of the \isa{axiomatization} block, which are literally the
equations the locale assumes.

\thy{enc-theorem}

\noindent
Read out, the theorem says that for all terminating $a$, $b$, $p_1$ and $p_2$,
it is not the case that $p_1$ codes a \BGA{} proof of $\enc{a = b}$ from no
hypotheses while $p_2$ codes a \BGA{} proof of $\enc{a \neq b}$ from no
hypotheses.  Every symbol in that statement, the quantifiers, the negation, the
conjunction and the proof checker, belongs to \QGA.
